\section{Overview}
\label{sec:overview}

Algorithm~7.5.1H, in Knuth's Pre-Fascicle~14A~\cite{Knuth}, is his
implementation of the Hopcroft--Karp algorithm~\cite{HK}. Inside
Section~7.5.1 Knuth calls it simply Algorithm~H, and so shall we. After a
greedy start (step~\Hs1)
it repeats a \emph{round}: build a dag that contains every shortest
augmenting path, then dismantle the dag, augmenting along a maximal
collection of vertex-disjoint shortest paths. Corollary~K bounds the
number of rounds by about $2\sqrt s$, where $s$ is the size of a maximum
matching, so the running time is $O\bigl((M+N+E)\sqrt s\,\bigr)$.

Is that square root ever needed? Knuth remarks on page~5 that graphs
needing anywhere near $\sqrt s$ rounds are hard to find, and that even
for such a graph ``its edges have to appear in a particularly perverse
order.''
Exercise~15 proposes a candidate: a \emph{Tchoukaillon graph}, on which
each round removes as few shortest paths as a simple counting argument
allows. The resulting round count is governed by the Tchoukaillon
numbers $\Ch_n$ of exercise~11, and since $\Ch_n\approx n^2/\pi$, such a
graph would force about $\sqrt{\pi s}$ rounds. Exercise~16 asks whether
Tchoukaillon graphs exist for every~$s$, and its answer raises an even
harder question: can some graph force \emph{more} rounds than a
Tchoukaillon graph would? Fig.~\ref{fig:rounds} shows the answer.

\begin{figure}[ht]
\centering
\begin{tikzpicture}
\begin{axis}[
  width=.92\linewidth, height=6.4cm,
  xmin=0.5, xmax=64.5, ymin=0, ymax=16,
  xtick={1,9,16,25,36,49,64}, ytick={2,4,...,16},
  xlabel={size $s$ of a maximum matching}, ylabel={rounds},
  axis lines=left, tick style={black!50}, tick label style={font=\scriptsize},
  label style={font=\small},
  ymajorgrids, grid style={black!8},
  legend style={draw=none,fill=white,fill opacity=.85,text opacity=1,
    at={(0.5,-0.24)},anchor=north,legend columns=2,font=\small,cells={anchor=west},/tikz/every even column/.append style={column sep=6mm}},
  clip=false]
\fill[sapcol!18] (axis cs:8.5,0) rectangle (axis cs:9.5,16);
\fill[sapcol!18] (axis cs:15.5,0) rectangle (axis cs:16.5,16);
\fill[sapcol!18] (axis cs:16.5,0) rectangle (axis cs:17.5,16);
\fill[sapcol!18] (axis cs:19.5,0) rectangle (axis cs:20.5,16);
\fill[sapcol!18] (axis cs:20.5,0) rectangle (axis cs:21.5,16);
\fill[sapcol!18] (axis cs:24.5,0) rectangle (axis cs:25.5,16);
\fill[sapcol!18] (axis cs:25.5,0) rectangle (axis cs:26.5,16);
\fill[sapcol!18] (axis cs:26.5,0) rectangle (axis cs:27.5,16);
\fill[sapcol!18] (axis cs:27.5,0) rectangle (axis cs:28.5,16);
\fill[sapcol!18] (axis cs:28.5,0) rectangle (axis cs:29.5,16);
\fill[sapcol!18] (axis cs:29.5,0) rectangle (axis cs:30.5,16);
\fill[sapcol!18] (axis cs:30.5,0) rectangle (axis cs:31.5,16);
\fill[sapcol!18] (axis cs:31.5,0) rectangle (axis cs:32.5,16);
\fill[sapcol!18] (axis cs:32.5,0) rectangle (axis cs:33.5,16);
\fill[sapcol!18] (axis cs:35.5,0) rectangle (axis cs:36.5,16);
\fill[sapcol!18] (axis cs:36.5,0) rectangle (axis cs:37.5,16);
\fill[sapcol!18] (axis cs:37.5,0) rectangle (axis cs:38.5,16);
\fill[sapcol!18] (axis cs:38.5,0) rectangle (axis cs:39.5,16);
\fill[sapcol!18] (axis cs:39.5,0) rectangle (axis cs:40.5,16);
\fill[sapcol!18] (axis cs:40.5,0) rectangle (axis cs:41.5,16);
\fill[sapcol!18] (axis cs:41.5,0) rectangle (axis cs:42.5,16);
\fill[sapcol!18] (axis cs:42.5,0) rectangle (axis cs:43.5,16);
\fill[sapcol!18] (axis cs:43.5,0) rectangle (axis cs:44.5,16);
\fill[sapcol!18] (axis cs:44.5,0) rectangle (axis cs:45.5,16);
\fill[sapcol!18] (axis cs:45.5,0) rectangle (axis cs:46.5,16);
\fill[sapcol!18] (axis cs:46.5,0) rectangle (axis cs:47.5,16);
\fill[sapcol!18] (axis cs:48.5,0) rectangle (axis cs:49.5,16);
\fill[sapcol!18] (axis cs:49.5,0) rectangle (axis cs:50.5,16);
\fill[sapcol!18] (axis cs:50.5,0) rectangle (axis cs:51.5,16);
\fill[sapcol!18] (axis cs:51.5,0) rectangle (axis cs:52.5,16);
\fill[sapcol!18] (axis cs:52.5,0) rectangle (axis cs:53.5,16);
\fill[sapcol!18] (axis cs:53.5,0) rectangle (axis cs:54.5,16);
\fill[sapcol!18] (axis cs:54.5,0) rectangle (axis cs:55.5,16);
\fill[sapcol!18] (axis cs:55.5,0) rectangle (axis cs:56.5,16);
\fill[sapcol!18] (axis cs:56.5,0) rectangle (axis cs:57.5,16);
\fill[sapcol!18] (axis cs:57.5,0) rectangle (axis cs:58.5,16);
\fill[sapcol!18] (axis cs:58.5,0) rectangle (axis cs:59.5,16);
\fill[sapcol!18] (axis cs:63.5,0) rectangle (axis cs:64.5,16);

\addplot[const plot mark mid,qcol,line width=1.3pt,mark=none,restrict x to domain=0:64.6]
  table[x=s,y=R]{figures/gen_rounds.dat};
\addlegendentry{$R(s)=\floor{\sqrt{4s+1}}-1$, the true maximum}
\addplot[const plot mark mid,chcol,line width=1.1pt,mark=none,densely dashed,restrict x to domain=0:64.6]
  table[x=s,y=T]{figures/gen_rounds.dat};
\addlegendentry{$T(s)$, a Tchoukaillon graph}
\node[font=\scriptsize,sapcol!80!black,anchor=south] at (axis cs:9,16) {$s=9$};
\end{axis}
\end{tikzpicture}

\caption{The largest possible number of rounds~$R(s)$ (blue), against
the number of rounds~$T(s)$ that a Tchoukaillon graph would need (red),
for matching sizes~$s\le64$. Shaded columns mark the sizes where
$R(s)>T(s)$; the first is $s=9$. The two staircases step up at the
quarter-squares $Q_n$ and at the Tchoukaillon numbers $\Ch_n$
(both defined below).}
\label{fig:rounds}
\end{figure}

\subsection*{The answers in brief}

\begin{summarybox}
\begin{enumerate}[leftmargin=1.4em,label=\textbf{\arabic*.}]
\item \textbf{The maximum number of rounds.} Count~\Hs1 as the first
round, as answer~15(c) does. A bipartite graph whose maximum matching
has $s$ edges can make Algorithm~H perform at most
\[
   R(s)=\bigfloor{\sqrt{4s+1}}-1
\]
rounds, whatever the order of its input, and for every $s$ some graph
and input order attain this bound (Theorem~R). Put differently,
$n$ rounds are possible exactly when
\[
   s\ \ge\ Q_n=\bigfloor{\frac{(n+1)^2}4}\qquad(\text{the quarter-squares }1,2,4,6,9,12,16,20,\ldots),
\]
while a hypothetical Tchoukaillon graph has at least $n$ rounds exactly
when $s\ge\Ch_n$ $(1,2,4,6,10,12,18,22,\ldots)$.
\item \textbf{Staircase graphs.} The extremal graphs are explicit, planar,
and have maximum degree three. Their input orders make Algorithm~H
find exactly one shortest augmenting path per round, at levels
$L=2,3,\ldots,n$, and we can prove every step of the execution
(Theorem~S). Algorithm~H performs $\Theta(E\sqrt V)$ elementary operations
on them (Theorem~W), so Corollary~K is tight even for sparse planar graphs.
\item \textbf{Beating Tchoukaillon.} Always $Q_n\le\Ch_n$, so the true
worst case is always at least as slow as a Tchoukaillon graph would be.
The first size where it is strictly slower is $s=9$: five rounds are
possible, against four. Asymptotically $R(s)-T(s)\sim(2-\sqrt\pi)\sqrt s$
(Section~\ref{sec:comparison}).
\item \textbf{The literal definition.} Exercise~15(c) asks that
\emph{every} augmenting path have the current shortest level at every
round. Graphs meeting that requirement exist only when
$s\in\{1,2,3,4,5,8,9\}$ (Corollary~L), because consecutive rounds of this
kind must at least halve the number of free girls (Theorem~U). Knuth's
16-pair example in answer~15(d) has augmenting paths of levels~2 and~3
immediately after~\Hs1, yet with suitable input orders it does follow
the Tchoukaillon schedule.
\item \textbf{The open question.} Can the Tchoukaillon \emph{schedule} (the
numbers of free girls and of paths removed, round by round) be realized
for every $s$? Explicit inputs exist for all $s\le47$, and they even keep
$f$ disjoint SAPs in every round, as Knuth's matrices do. Any family of
inputs that realizes the schedule for unboundedly large $s$ must contain
at least $s/12-O(1)$ independent cycles (Theorem~C). Some candidate graphs
fail only because of the order in which Algorithm~H searches: its queue,
the head insertion of step~\Hs6, and the replacement rule of step~\Hs{16}
(Section~\ref{sec:order}).
\end{enumerate}
\end{summarybox}

\subsection*{How these notes are organized}

Section~\ref{sec:background} recalls just enough of Algorithm~H to make
precise what an ``input'' is and which choices it leaves open; it
also works through one round of an example that recurs throughout.
Section~\ref{sec:ex15} restates exercise~15, explains where the
Tchoukaillon numbers come from, and separates three ways of reading the
definition of a Tchoukaillon graph. Sections~\ref{sec:upper}
and~\ref{sec:staircase} prove Theorem~R: first the upper bound, which is
a sharpened form of the counting in Theorem~H, then the staircase
construction. Section~\ref{sec:comparison} compares the answer with the
Tchoukaillon numbers. Section~\ref{sec:literal} settles the literal
reading of exercise~15, and Section~\ref{sec:operational} collects what
we know about the open, numerical reading. Section~\ref{sec:sparse}
gives simpler examples for the running-time question alone.
Section~\ref{sec:credits} discusses the literature.

The appendices give a complete ordered input for Knuth's 16-pair graph,
details of a family that meets the literal definition, a catalog of
computed examples, and a guide to the accompanying programs.

\paragraph{Which draft.} Page, exercise, step, and equation numbers refer
to the zeroth printing of Pre-Fascicle~14A, revision~$-88$, dated
18~September 2026. A pre-fascicle is a working draft, so these numbers may
change in later revisions.

\subsection*{What is classical and what is new}

The tools are the ones in the pre-fascicle: Berge's augmenting paths
\cite{Berge}, the layered dag and the fact that the shortest level
increases from round to round~\cite{HK}, and the counting argument
behind Theorem~H. The upper bound in Theorem~R refines Hopcroft and Karp's
bound on the number of phases. As far as we know, the staircase graphs,
the exact value of $R(s)$, and the results of Sections~\ref{sec:literal}
and~\ref{sec:operational}, are new; Section~\ref{sec:credits} discusses
the literature.
